% =====================================================================
%  Thesis proposal, Appendix: Coupled Entry Dynamics and Polynomial Embedding
%  Owns: apx:dyn
%
%  Transcription of notes/embedding.md at commit 18cc0a2 (the author's own
%  text) into LaTeX. Wording follows the note; nothing is added, removed, or
%  strengthened. Two spelling fixes are marked in comments below.
% =====================================================================

\section{Coupled Entry Dynamics and Polynomial Embedding}
\label{apx:dyn}

\subsection{Coupled 6-DOF Entry Dynamics}
\label{sub:apx_dyn_6dof}

The entry vehicle state comprises the coupled translational, rotational,
thermal, ablative, and modal subsystems. Position is carried in Cartesian
Earth-centred Earth-fixed (ECEF) components and attitude as a quaternion from
the ECEF frame to the body frame. No angle is a state, so the dynamics have no
coordinate singularities anywhere on the globe; geocentric latitude and
longitude are read out from the position after integration.

\paragraph{Frames.}
\begin{itemize}
  \item $\mathcal E$: Earth-centred Earth-fixed (ECEF), rotating with the
        constant planetary rate $\vec\omega_E = \omega_E\,\hat z_{\mathcal E}$.
  \item $\mathcal B$: body frame, fixed at the instantaneous center of mass.
\end{itemize}

The attitude quaternion $\vec q = (q_0, q_1, q_2, q_3)^\top$ carries
$\mathcal E \to \mathcal B$, with direction-cosine matrix
\begin{equation}
\label{eq:dcm}
\mathbf R \equiv \mathbf R_{B/E}(\vec q)=
\begin{pmatrix}
q_0^2+q_1^2-q_2^2-q_3^2 & 2(q_1q_2+q_0q_3) & 2(q_1q_3-q_0q_2)\\
2(q_1q_2-q_0q_3) & q_0^2-q_1^2+q_2^2-q_3^2 & 2(q_2q_3+q_0q_1)\\
2(q_1q_3+q_0q_2) & 2(q_2q_3-q_0q_1) & q_0^2-q_1^2-q_2^2+q_3^2
\end{pmatrix},
\end{equation}
quadratic in $\vec q$.

\paragraph{Position Kinematics.}
With $\vec r = (r_1, r_2, r_3)^\top$ the ECEF position and
$\vec V = (u, v, w)^\top$ the planet-relative velocity resolved in
$\mathcal B$, the kinematic equation is:
\begin{equation}
\label{eq:poskin}
\dot{\vec r} = \mathbf R^\top \vec V.
\end{equation}

\paragraph{Translational Dynamics.}
Newton's law holds in the inertial frame; transporting the planet-relative
velocity through $\mathcal E$ into $\mathcal B$, with inertial body rates
$\vec\omega = (p, q, r_b)^\top$, mass $m$, aerodynamic force $\vec F_a$, thrust
$\vec F_t$, and the planetary rate resolved in the body frame
$\vec\omega_E^{\mathcal B} = \mathbf R\,\omega_E\hat z_{\mathcal E}$, the
translational dynamical equations are:
\begin{equation}
\label{eq:transdyn}
\dot{\vec{V}}\big|_{\mathcal{B}} = \frac{1}{m}(\vec{F}_a + \vec{F}_t)
+ \mathbf R\,\vec{g}(\vec r)
- (\vec\omega + \vec\omega_E^{\mathcal B})\times\vec{V}
- \vec\omega_E^{\mathcal B}\times(\vec\omega_E^{\mathcal B}\times\mathbf R\,\vec{r}).
\end{equation}

Gravity, including the $J_2$ oblateness term, in ECEF components:
\begin{equation}
\label{eq:gravity}
\vec g(\vec r) = -\frac{\mu}{\|\vec r\|^3}\left[\left(1
+ \frac{3}{2}J_2\frac{R_e^2}{\|\vec r\|^2}\left(1
- 5\frac{r_3^2}{\|\vec r\|^2}\right)\right)\vec r
+ 3J_2\frac{R_e^2}{\|\vec r\|^2}\,r_3\,\hat z_{\mathcal E}\right].
\end{equation}

\begin{remark}
This is the gradient of the $J_2$ potential
$\frac{\mu}{\|\vec r\|}\bigl[1 - \tfrac12 J_2 (R_e/\|\vec r\|)^2(3 r_3^2/\|\vec r\|^2 - 1)\bigr]$;
its radial and northward components reproduce the spherical form, with
$r_3/\|\vec r\| = \sin\phi$.
\end{remark}

\paragraph{Rotational Dynamics / Euler's Equations.}
With angular momentum $\vec h = \mathbf I\vec\omega$, combined aerodynamic /
reaction control system (RCS) moments
$\vec M = \vec M_a + \vec M_{\mathrm{rcs}}$, and the inertia tensor
$\mathbf I = \mathbf I(s, \eta)$ depending on recession depth $s$ and modal
amplitudes $\eta$ (so that $\dot{\mathbf I}\,\vec\omega$ is the variable-inertia
term), the rotational dynamical equations are:
\begin{equation}
\label{eq:rotdyn}
\mathbf{I}\,\dot{\vec\omega} = \vec{M} - \vec\omega\times(\mathbf{I}\vec\omega)
- \dot{\mathbf{I}}\,\vec\omega.
\end{equation}

\paragraph{Quaternion Kinematics.}
The quaternion is driven by the body rate \emph{relative to} $\mathcal E$,
$\vec\omega_{B/E} = \vec\omega - \vec\omega_E^{\mathcal B}$. For
$\vec a = (a_1, a_2, a_3)^\top$,
\begin{equation}
\label{eq:qkin}
\dot{\vec{q}} = \tfrac{1}{2}\,\Omega(\vec\omega
- \mathbf R\,\omega_E\hat z_{\mathcal E})\,\vec{q},
\qquad
\Omega(\vec a) = \begin{pmatrix}
0 & -a_1 & -a_2 & -a_3\\
a_1 & 0 & a_3 & -a_2 \\
a_2 & -a_3 & 0 & a_1 \\
a_3 & a_2 & -a_1 & 0
\end{pmatrix}.
\end{equation}

\begin{remark}
The rate is \textbf{polynomial} --- bilinear in $(\vec\omega, \vec q)$ and cubic
in $\vec q$ through the Earth-rate term. With this $\Omega$ and $\mathbf R$ of
\cref{eq:dcm}, $\dot{\mathbf R} = -[\vec\omega_{B/E}\times]\,\mathbf R$ holds
identically, and $\Omega(\vec a)$ is skew for every $\vec a$, so $|\vec q|$ is
preserved.
\end{remark}

\paragraph{Thermal.}
% note: "sqaure-root" in the note, spelling corrected here
With the density-square-root embedding $y = \sqrt{\rho/\rho_0}$ and
$V = \sqrt{u^2+v^2+w^2}$, the modulus lift $\dot\chi = c\,\chi\,\dot T_w$
replaces $T_w$ with a polynomial via Varma--Savageau--Voit recasting.
Stagnation-point heating (Sutton--Graves, with nose radius $R_n$, sea-level
density $\rho_0$) equations are then given by:
\begin{equation}
\label{eq:thermal}
\dot Q_{\mathrm{tot}} = \dot Q = k_Q\sqrt{\rho_0/R_n}\,y\,V^3,
\qquad
\dot T_w = \frac{\dot Q - (T_w - T_\infty)/\tau_\text{cond}}{C_w}.
\end{equation}

\paragraph{Ablation.}
With recession rate coupled to heat flux and mass loss from char-layer
blow-off, the ablative equations are given by:
\begin{equation}
\label{eq:ablation}
\dot s = \alpha_s \,\dot Q, \qquad \dot m = -\dot m_\text{abl}.
\end{equation}

\paragraph{Modal/Structural Elasticity.}
With piston-theory forcing and temperature-dependent natural frequency, per
mode $i$ our modal equations are:
\begin{equation}
\label{eq:modal}
\ddot\eta_i + 2\zeta_i\,\omega_i(T)\,\dot\eta_i + \omega_i(T)^2\,\eta_i
= Q_i(x, \eta, \dot\eta).
\end{equation}

\subsection{Polynomial Embedding via Adjunction}
\label{sub:apx_dyn_embedding}

A priori, the dynamics contain transcendental and rational terms. We repair
each of these by adjoining one variable and one polynomial defining constraint,
so the entire augmented system is polynomial in the augmented state:

\begin{enumerate}
  \item \textbf{Transcendental Incidence:} $\alpha = \operatorname{atan2}(w, u)$,
        $\beta = \arcsin(v/V)$ are inverse-trigonometric functions of body
        rates, we replace them with polynomial direction-cosine ratios
        $(u/V, v/V, w/V)$.
  \item \textbf{Square-Root Speed:} Since $V = \sqrt{u^2+v^2+w^2}$ carries a
        square root, we adjoin $V$ with the velocity magnitude constraint
        $V^2 = u^2+v^2+w^2$.
  \item \textbf{Exponential Atmosphere:} Since $\rho = \rho_0 e^{-h/H_s}$ with
        $h = \|\vec r\| - R_e$ appears in $\bar q = \tfrac12\rho V^2$, we adjoin
        $y = \sqrt{\rho/\rho_0}$ with rate
        $\dot y = -\frac{y}{2H_s}\,\tfrac{d}{dt}\|\vec r\|$, where the radial
        rate $\tfrac{d}{dt}\|\vec r\| = \varrho\,(\mathbf R\vec r)\cdot\vec V$ is
        polynomial once $\varrho$ (item 4) is adjoined.
  \item \textbf{Reciprocal Radius:} Gravity carries $\|\vec r\|^{-3}$,
        $\|\vec r\|^{-5}$, $\|\vec r\|^{-7}$, so we adjoin
        $\varrho = 1/\|\vec r\|$ with $\varrho^2\,\|\vec r\|^2 = 1$,
        $\varrho > 0$, and rate
        $\dot\varrho = -\varrho^3\,(\mathbf R\vec r)\cdot\vec V$. Gravity is then
        polynomial:
        \begin{equation}
        \label{eq:gravity_poly}
        \vec g = -\mu\varrho^3\left[\left(1
        + \tfrac32 J_2 R_e^2\varrho^2\left(1 - 5 r_3^2\varrho^2\right)\right)\vec r
        + 3J_2R_e^2\varrho^2\,r_3\,\hat z_{\mathcal E}\right].
        \end{equation}
        No trigonometric adjunctions are needed: no angle is a state, and
        latitude dependence enters only through $r_3\varrho = \sin\phi$.
  \item \textbf{Quaternion Norm:} $|\vec q|^2 = 1$ is already a polynomial
        invariant (skew-symmetry of $\Omega$ preserves it), no adjunction
        necessary.
  \item \textbf{Heat / Load-Factor Square Roots:} $\dot Q \propto y V^3$ is
        already polynomial; however, $n = \|\vec F_a\|/(mg_0)$ requires squaring
        to $(\bar q S)^2(C_X^2+C_Y^2+C_Z^2) \le (n_\text{max} m g_0)^2$.
\end{enumerate}

\paragraph{The Augmented State.}
Collecting all coordinates:
\begin{equation}
\label{eq:augmented_state}
\tilde{x} = \bigl(
\underbrace{r_1,\,r_2,\,r_3,\,u,\,v,\,w,\,q_0,\,q_1,\,q_2,\,q_3,\,p,\,q,\,r_b}_{13,\ \text{core}},\;
\underbrace{y,\,V,\,\varrho,\,Q_{\mathrm{tot}},\,m,\,s,\,\chi}_{7,\ \text{adjoined}},\,
\underbrace{\eta,\,\dot\eta}_{2n_\eta,\ \text{modal}}\bigr).
\end{equation}
